\chapter{Generalidades y datos descriptivos}
\label{cap:generalidades}

\opening{Un manual que no está al día no sirve para operar.}

\section{Introducción al manual de operación}

El Manual de Operación del Piloto (\emph{Pilot's Operating Handbook}, POH) es la guía de operación del R22. Contiene el material que los reglamentos 14 CFR partes 21 y 27 exigen entregar al piloto y datos complementarios del fabricante. Sus secciones 2, 3, 4 y 5 están aprobadas por la FAA y constituyen el \emph{Rotorcraft Flight Manual} que debe llevarse a bordo en todo momento.\fuente{POH, sección 1}

El POH no sustituye una instrucción de vuelo adecuada y competente ni el conocimiento de las directivas de aeronavegabilidad vigentes, los reglamentos y las circulares de asesoramiento. Tampoco es una guía de instrucción básica ni un manual de entrenamiento, y no debe usarse para fines operacionales si no se mantiene actualizado.

Asegurar que el helicóptero esté en condición aeronavegable es responsabilidad del propietario. El piloto al mando es responsable de determinar que el helicóptero está seguro para el vuelo y de mantenerse dentro de las limitaciones de operación indicadas por las marcas de los instrumentos, los letreros (\emph{placards}) y el manual.

Como es muy difícil consultar un manual mientras se vuela un helicóptero, el piloto debe estudiar el manual completo y familiarizarse a fondo con las limitaciones, la performance, los procedimientos y las características de manejo del helicóptero \emph{antes} de volar.

El POH está dividido en diez secciones numeradas. Las limitaciones y los procedimientos de emergencia van antes que los procedimientos normales y la performance para que sean más fáciles de encontrar. Algunos números de párrafo, figuras y páginas se omiten deliberadamente (``páginas intencionalmente en blanco'') para permitir futuras ampliaciones.

\begin{nota}
Suscripción a actualizaciones: Robinson ofrece una suscripción ``Clase A'' que envía los cambios del POH y las nuevas Safety Notices (45 USD por dos años, en \texttt{www.robinsonheli.com}). La fecha estampada en cada manual refleja la revisión con la que fue armado; consulte el sitio web para conocer la revisión más reciente.
\end{nota}

\section{Avisos y notas}

Las \textbf{precauciones} (\emph{CAUTION}) indican que pueden producirse daños al equipo, lesiones o la muerte si no se sigue el procedimiento o la instrucción. Las \textbf{notas} (\emph{NOTE}) dan énfasis o información complementaria. En el Manual de Mantenimiento se agrega la \textbf{advertencia} (\emph{WARNING}): riesgo de lesión personal o muerte.

\section{Dimensiones exteriores}

\widefigure[0.82\linewidth]{assets/poh/dimensiones.png}{Dimensiones exteriores del R22 (aeronave de producción tardía). Radio del rotor principal 151~in (3,84~m); diámetro del rotor de cola 42~in (1,07~m); longitud del fuselaje y cono de cola 259~in (6,58~m); longitud total 345~in (8,76~m); altura del techo de cabina 69~in (1,75~m); altura total 107~in (2,72~m); ancho de cabina 44~in (1,12~m); trocha de patines 76~in (1,93~m).}

\section{Datos descriptivos}

\subsection{Rotor principal}
\begin{checklist}
\chk{Articulación}{Libre para balancear y conificar; rígido en el plano}
\chk{Número de palas}{2}
\chk{Diámetro}{25 ft 2 in (7,67 m)}
\chk{Cuerda de la pala}{7,2 in interior; 7,7 in exterior}
\chk{Torsión de la pala}{$-8^\circ$}
\chk{Velocidad de punta a 104\,\% RPM}{698 ft/s}
\end{checklist}

\subsection{Rotor de cola}
\begin{checklist}
\chk{Articulación}{Libre para balancear; rígido en el plano}
\chk{Número de palas}{2}
\chk{Diámetro}{42 in (1,07 m)}
\chk{Cuerda de la pala}{4 in (constante)}
\chk{Torsión de la pala}{$0^\circ$}
\chk{Ángulo de preconicidad}{1° 11'}
\chk{Velocidad de punta a 104\,\% RPM}{622 ft/s}
\end{checklist}

\subsection{Sistema de transmisión}
\begin{checklist}
\chk{Motor a polea superior}{Dos correas en V dobles; reducción 0,8536:1}
\chk{Polea superior a eje de transmisión}{Embrague de rueda libre tipo \emph{sprag}}
\chk{Eje a rotor principal}{Engranajes cónicos espirales; reducción 11:47}
\chk{Eje a rotor de cola}{Engranajes cónicos espirales; multiplicación 3:2}
\end{checklist}

\subsection{Grupo motor}

Un motor Lycoming O-320 u O-360 de cuatro cilindros horizontalmente opuestos, de accionamiento directo, enfriado por aire, con carburador y de aspiración normal. Cilindrada 319,8~in³ (O-320) o 361,0~in³ (O-360). El enfriamiento lo provee un ventilador tipo jaula de ardilla de accionamiento directo.

\begin{table}[htbp]
\centering\small
\caption{Motores y potencias del R22.}
\label{tab:motores}
\begin{tabularx}{\linewidth}{@{}L{2.0cm}Y@{}}
\toprule
Motor & Potencia nominal y aplicación \\
\midrule
O-320-A2B o -A2C & 150 BHP a 2700 RPM (R22 Standard) \\
O-320-B2C & 160 BHP a 2700 RPM (R22 HP, Alpha y Beta) \\
O-360-J2A & 145 BHP (reducida) a 2700 RPM (R22 Beta II) \\
\midrule
\multicolumn{2}{@{}l}{\textbf{Potencias autorizadas en el R22}} \\
MCP & 124 BHP a 2652 RPM (104\,\% en el tacómetro) \\
TOP (5 min) & 131 BHP a 2652 RPM, solo Beta y Beta II \\
\bottomrule
\end{tabularx}
\end{table}

Los grados y la capacidad de combustible están en el capítulo~\ref{cap:limitaciones}; los grados y la capacidad de aceite, en el capítulo~\ref{cap:manejo}.

\section{Definiciones de performance}

\begin{longtable}{@{}L{1.9cm}p{\dimexpr\linewidth-1.9cm-2\tabcolsep}@{}}
\toprule
Término & Significado \\
\midrule
\endhead
KIAS & Velocidad indicada en nudos: la que muestra el anemómetro. \\
KCAS & Velocidad calibrada en nudos: la indicada corregida por error de instrumento y de posición (ver la curva de calibración en el capítulo~\ref{cap:performance}). \\
KTAS & Velocidad verdadera en nudos: velocidad respecto al aire no perturbado; es la KCAS corregida por altitud de presión y temperatura. \\
\textsc{v}\textsubscript{NE} & Velocidad que nunca debe excederse. \\
\textsc{v}\textsubscript{Y} & Velocidad de máxima razón de ascenso. \\
Altitud MSL & Altitud sobre el nivel del mar indicada por el altímetro (corregido por errores de posición e instrumento) con la subescala barométrica ajustada a la presión existente al nivel del mar. \\
Altitud de presión & Altitud indicada por el altímetro (corregido) con la subescala ajustada a 29,92~inHg (1013,2~mb). \\
Altitud de densidad & Altitud, en condiciones ISA, a la que el aire tendría la misma densidad; es la altitud de presión corregida por la temperatura exterior. \\
ISA & Atmósfera estándar internacional: 29,92~inHg y 15~°C al nivel del mar, con una disminución de temperatura de 1,98~°C cada 1000~ft. \\
BHP & Potencia al freno: potencia real entregada por el motor. \\
MAP & Presión de admisión: presión absoluta, en pulgadas de mercurio, en el múltiple de admisión del motor. \\
RPM & Revoluciones por minuto del motor o del rotor principal. El tacómetro las muestra como porcentaje de 2550~RPM de motor o de 510~RPM de rotor. \\
MCP & Potencia máxima continua. \\
TOP & Potencia de despegue (normalmente por un máximo de 5 minutos). \\
Altitud crítica & Altitud a la que el acelerador completamente abierto produce la máxima potencia permitida (MCP o TOP). \\
TOGW & Peso bruto de despegue. \\
OAT & Temperatura del aire exterior. \\
CAT & Temperatura del aire del carburador. \\
CHT & Temperatura de culata de cilindro. \\
GPH & Galones por hora. \\
AGL & Sobre el nivel del terreno. \\
IGE & Con efecto suelo. \\
OGE & Fuera del efecto suelo. \\
ALT & Alternador. \\
\bottomrule
\end{longtable}

\section{Definiciones de peso y balance}

\begin{longtable}{@{}L{2.6cm}p{\dimexpr\linewidth-2.6cm-2\tabcolsep}@{}}
\toprule
Término & Significado \\
\midrule
\endhead
Dato de referencia & Plano vertical desde el que se miden las distancias horizontales para el balance. En el R22, el dato longitudinal está 100~in por delante del eje del mástil del rotor principal. \\
Estación & Ubicación longitudinal a lo largo del fuselaje, expresada en pulgadas desde el dato de referencia. \\
Brazo & Distancia horizontal desde el dato de referencia hasta el centro de gravedad (CG) de un elemento. \\
Momento & Peso del elemento multiplicado por su brazo. \\
Centro de gravedad (CG) & Ubicación a lo largo del fuselaje (normalmente en pulgadas desde el dato) en la que el helicóptero quedaría en equilibrio. Se calcula dividiendo el momento total por el peso total. \\
Límites de CG & Ubicaciones extremas del CG dentro de las que debe operarse el helicóptero para un peso dado. \\
Combustible utilizable & Combustible disponible para la planificación del vuelo. \\
Combustible no utilizable & Combustible que queda en el tanque y que no puede proporcionar un flujo ininterrumpido de manera confiable en la actitud de vuelo crítica. \\
Peso vacío estándar & Peso del helicóptero estándar incluyendo el combustible no utilizable, todos los fluidos operativos y el aceite de motor completo. \\
Peso vacío básico & Peso vacío estándar más el peso del equipamiento opcional instalado. \\
Carga paga & Peso de ocupantes, carga y equipaje. \\
Carga útil & Diferencia entre el peso bruto máximo y el peso vacío básico. \\
\bottomrule
\end{longtable}
